\section{Task Formulation}
\label{sec:problem}

An agent accumulates an interaction history $H=(s_1,\ldots,s_T)$ across
dated dialogue sessions. A memory construction process transforms this
history into a memory $\mem$ containing source passages and derived
representations. Given a question $q$, the retrieval task is to construct
a context $\ctx$ within a budget $b$, from which a reader produces an
answer:
\begin{equation}
 \ctx=\mathrm{Retrieve}(q,\mem;b),\qquad
 \hat a=\mathrm{Read}(q,\ctx).
 \label{eq:task}
\end{equation}
The budget limits the amount of evidence delivered to the reader.
The objective is to select evidence that supports an accurate answer,
rather than merely collecting records that resemble the question.

We focus on questions whose relational requirements can be represented
by a conjunctive plan. A plan describes the facts to be found and the
variables through which they must connect. Retrieval then seeks
witnesses of that plan in the stored memory
(\cref{sec:preliminaries}). Temporal information and the requested answer
form further guide the search and the organization of the resulting
evidence.

This formulation separates two responsibilities: interpreting the
question as a retrieval plan and finding memory records that satisfy
its relational requirements. Witnesses make the latter explicit.
Their relationship to the original question depends on the quality
of the interpretation and stored facts, as discussed in
\cref{sec:discussion}.
